\title{Nominal Budgets and Realized Resources in Closed-API Large-Language-Model Inference: A Performance Evaluation Protocol}

\author[aff1]{Soroush Vahidi\corref{cor1}\fnref{fn1}}
\ead{sv96@njit.edu}

\cortext[cor1]{Corresponding author}
\fntext[fn1]{Sole author. ORCID: \href{https://orcid.org/0000-0003-1934-6282}{0000-0003-1934-6282}.}

\address[aff1]{Ying Wu College of Computing, New Jersey Institute of Technology, Newark, New Jersey, USA}

\begin{abstract}
Nominal inference budgets are convenient for large language model (LLM) systems, but equal logical
entitlement need not imply equal successful completions, retries, tokens, latency, cost, or selector
attribution. We present a protocol for closed-API, budgeted inference that separates discovery from
commit-rule selection, evaluates identical-pool controls, records a componentwise realized resource
vector, labels provider-transfer evidence, and preserves protocol-blocked outcomes. On a 15-cell
provider-by-dataset matrix ($n{=}3394$; Fireworks$\times$GPQA-Diamond protocol-blocked),
identical-pool plurality over four generators (Pooled-4) exceeds a failure-trace gated selector
(FTA): $66.53\%$ versus $65.00\%$ (McNemar $p{=}0.00027$). The dominant FTA gate reverses sign under
provider transfer, including $2/18$ and $7/31$ rescue/regression patterns on Azure mathematics
cells. A resource audit changes efficiency interpretation without changing accuracy: reconstructed
Frontier successful completions rise from $2.78$ to $5.38$ of nominal budget $B{=}6$, while token/USD
fields remain lower bounds. The contribution is the protocol: nominal budgets alone are insufficient
for defensible closed-API inference-performance conclusions.
\end{abstract}

\begin{keyword}
budgeted inference \sep performance measurement \sep resource accounting \sep closed-API inference \sep large language models
\end{keyword}
